\documentclass[12pt]{article}
\usepackage{sbc-template}
\usepackage[english,brazilian]{babel}
\usepackage[utf8]{inputenc}
\usepackage[T1]{fontenc}
\usepackage{graphicx}
\usepackage{url}
\usepackage{booktabs}
\usepackage{amsmath}
\usepackage{amssymb}
\usepackage{tikz}
\usetikzlibrary{arrows.meta, positioning, calc, fit, backgrounds}
\usepackage{microtype}
\usepackage[alf,abnt-etal-list=0]{abntex2cite}

\title{SIM --- Sistema de Informações Municipais:\\
Acesso Conversacional a Dados Abertos Governamentais}

\author{Rafael A. S. Rezende\inst{1}, Denilson A. Pereira\inst{1}}

\address{Departamento de Ciência da Computação -- Universidade Federal de Lavras (UFLA)\\
  Caixa Postal 3037 -- 37200-900 -- Lavras -- MG -- Brasil\\
  \email{rafael.rezende@estudante.ufla.br, denilsonpereira@ufla.br}
}

\begin{document}

\maketitle

\begin{otherlanguage}{english}
\begin{abstract}
Municipal open data is essential for social accountability, yet effective public access remains constrained by complex transparency portals and relational database schemas. We present SIM (Sistema de Informações Municipais), a multi-instance conversational platform that integrates document retrieval over official municipal gazettes and analytical query generation via Text-to-SQL over state accounting (SICOM/TCE-MG) and federal (CGU) databases. We evaluate the pilot instance of Caeté, Minas Gerais, deployed over WhatsApp with LLM-based orchestration. Empirical validation was conducted across three experimental batches: Text-to-SQL on 80 executable cases derived from ten analytical intents (Baseline E1: 97.5\% VSR and 62.5\% EF; domain-adapted Prompt B: 100\% VSR and 72.5\% EF, 95\% CI for EF: 61.9--81.1\%) and document retrieval over 40 stratified queries (100\% Recall@5 and 97.5\% MRR under hybrid search versus 22.5\% under dense search alone, demonstrating the necessity of hybrid retrieval for numbered legal acts). All experimental artifacts are publicly available for offline replication.
\end{abstract}
\end{otherlanguage}

\begin{resumo}
Dados abertos municipais são instrumentos essenciais de fiscalização social, mas o acesso efetivo permanece restrito a usuários com domínio técnico de portais de transparência e linguagens de consulta estruturada. Este artigo apresenta o SIM (Sistema de Informações Municipais), plataforma conversacional multi-instância que integra recuperação documental sobre o Diário Oficial e consulta analítica via Text-to-SQL sobre bases contábeis estaduais do SICOM e federais da Controladoria-Geral da União. Avalia-se a instância piloto do município de Caeté, em Minas Gerais, integrada ao WhatsApp com orquestração por modelos de linguagem. A validação empírica organizou-se em três lotes experimentais: Text-to-SQL com 80 casos executáveis derivados de dez intenções analíticas centrais (linha de base E1 com VSR de 97,5\% e EF de 62,5\%; Prompt B com adaptação de domínio com VSR de 100\% e EF de 72,5\%, IC 95\% da EF: 61,9--81,1\%) e recuperação documental com 40 perguntas estratificadas (Recall@5 de 100\% e MRR de 97,5\% na configuração híbrida, ante 22,5\% na busca densa isolada, comprovando a necessidade de recuperação híbrida para atos normativos numerados). Os artefatos experimentais estão disponíveis publicamente para replicação offline.
\end{resumo}

\section{Introdução}

A Lei de Acesso à Informação (Lei nº 12.527/2011) consolidou o direito de consulta a dados governamentais no Brasil. Em nível municipal, sistemas como o Sistema Informatizado de Contas dos Municípios (SICOM)\footnote{SICOM/TCE-MG: \url{https://dadosabertos.tce.mg.gov.br/}}, mantido pelo Tribunal de Contas do Estado de Minas Gerais (TCE-MG), e portais federais complementares, a exemplo da Controladoria-Geral da União (CGU)\footnote{Portal da Transparência: \url{https://portaldatransparencia.gov.br/}} e da Secretaria do Tesouro Nacional por meio do Siconfi\footnote{Siconfi: \url{https://www.tesourotransparente.gov.br/consultas/consultas-siconfi/}}, reúnem registros contábeis de despesas, receitas, contratos e folhas de pagamento. Contudo, a consulta efetiva permanece restrita a especialistas capazes de navegar em esquemas relacionais complexos e vocabulários fiscais controlados.

Nesse cenário, os modelos de linguagem de larga escala (LLMs) viabilizam novas formas de interação entre cidadãos e bases governamentais. Por um lado, as técnicas de tradução de linguagem natural para SQL (Text-to-SQL) convertem perguntas livres diretamente em consultas estruturadas executáveis sobre bancos de dados relacionais \cite{Yu2018Spider,Li2024BIRD,Qu2025CodeS}. Por outro lado, a recuperação aumentada por geração (RAG) permite consultar acervos documentais não estruturados, recuperando trechos textuais relevantes para fundamentar a resposta gerada pelo modelo \cite{Lewis2020RAG,Gao2024RAGSurvey,Sharma2025RAGSurvey}. A integração dessas duas abordagens no contexto da tecnologia cívica (GovTech) oferece uma oportunidade concreta para reduzir barreiras técnicas e democratizar o acesso à transparência pública \cite{Janssen2020GovTech,Lei2024SICOM}.

Apesar desse potencial, a aplicação prática dessas tecnologias a dados municipais brasileiros impõe desafios relevantes. Destacam-se a heterogeneidade e a extensão dos esquemas contábeis, a presença frequente de dados incompletos ou ruidosos, a coexistência de dados financeiros tabulares e atos normativos publicados no Diário Oficial, além da exigência de respostas com baixa latência em canais de comunicação acessíveis, como o WhatsApp.

Diante desse panorama, este trabalho investiga a seguinte questão de pesquisa: \emph{É viável integrar recuperação documental e Text-to-SQL em canal conversacional para consulta a dados públicos municipais, e qual a eficácia de regras de domínio no \textit{prompt} --- o texto de instrução em linguagem natural fornecido ao modelo antes de cada consulta, contendo o esquema do banco de dados e as regras de negócio aplicáveis --- frente ao uso de modelos de linguagem de maior porte na fidelidade das respostas tabulares?}

Como hipótese de trabalho, assume-se que: (i)~a orquestração entre RAG e Text-to-SQL é funcionalmente viável para atender demandas cidadãs sobre dados municipais abertos, e (ii)~a inclusão de regras específicas do domínio contábil no contexto do modelo proporciona maior fidelidade funcional do que a simples substituição do modelo por uma versão de maior escala, mantendo a latência e os custos operacionais em níveis compatíveis com canais síncronos de mensageria.

Para investigar essas hipóteses, desenvolveu-se o SIM (Sistema de Informações Municipais), uma plataforma conversacional concebida como uma arquitetura genérica e replicável para municípios brasileiros. A escolha do município de Caeté (MG) (código municipal 3110004 do Instituto Brasileiro de Geografia e Estatística, IBGE) para a instância piloto decorreu da proximidade a gestores e atores públicos municipais para levantamento de demandas reais de transparência cívica, combinada ao interesse mútuo de viabilizar a fiscalização cidadã. Ademais, a base contábil de Caeté reúne dados de 2020 a 2025 com lacunas históricas reais de publicação em folha de pagamento e servidores de 2020 a 2022, o que permitiu testar a resiliência do sistema sob condições autênticas do setor público.

As contribuições principais deste trabalho são:
\begin{enumerate}
    \item Arquitetura do SIM, integrando orquestração cognitiva por modelo de linguagem, recuperação documental e tradução Text-to-SQL em canal conversacional no WhatsApp;
    \item Pipeline reprodutível de extração, transformação e carga (ETL) para normalização e exposição analítica de dados contábeis e normativos municipais;
    \item Conjunto de referência estruturado com 80 casos de consulta SQL e 40 consultas documentais para avaliação experimental;
    \item Protocolo de avaliação quantitativo baseado em taxa de SQL válida, equivalência funcional com tolerância fiscal, métricas de recuperação documental (Recall@$k$, presença do trecho correto entre os $k$ primeiros resultados, e MRR, do inglês \textit{Mean Reciprocal Rank}) e análise de significância estatística --- todas formalmente definidas na Seção 4.2;
    \item Avaliação empírica comparativa evidenciando o impacto de regras de domínio na geração de SQL e a necessidade de recuperação híbrida para atos normativos numerados.
\end{enumerate}

\section{Trabalhos Relacionados}

Esta seção analisa os principais trabalhos relacionados na literatura, organizados em três eixos: interfaces Text-to-SQL e seus benchmarks (Seção 2.1), recuperação documental em textos governamentais (Seção 2.2) e aplicações de GovTech para transparência pública (Seção 2.3).

\subsection{Interfaces em Linguagem Natural e Benchmarks Text-to-SQL}

O avanço de interfaces em linguagem natural para bancos de dados foi impulsionado por conjuntos de referência públicos. O Spider \cite{Yu2018Spider} estabeleceu a acurácia de execução (Execution Accuracy, EX) como métrica padrão em esquemas sintéticos em inglês. O BIRD \cite{Li2024BIRD} ampliou o desafio para bancos de dados volumosos com valores ruidosos e conhecimento externo, demonstrando que modelos comerciais de linguagem (como o GPT-4) frequentemente obtêm EX inferior a 60\% em condições realistas \cite{Li2024BIRD,Pourreza2024DIN,Zheng2026GovLLM}.

Métodos recentes exploram decomposição de consultas, seleção de exemplos e vinculação de esquema (\textit{schema linking}). O DIN-SQL \cite{Pourreza2024DIN} decompõe a geração em etapas sequenciais com autocorreção; o DAIL-SQL \cite{Gao2024DAILSQL} otimiza a seleção de exemplos em contexto; o RAT-SQL \cite{Wang2020RATSQL} codifica relações estruturais do esquema, e frameworks como CHESS \cite{Talaei2024CHESS} e CodeS \cite{Qu2025CodeS} estruturam a síntese modular de SQL. Diferentemente desses métodos que operam com múltiplos turnos de modelo e elevado consumo de tokens, o SIM prioriza baixa latência em requisição única, adaptando o contexto com regras de domínio compiladas. Levantamentos recentes consolidam esse cenário fragmentado de técnicas \cite{Hong2025Survey}, ao passo que trabalhos de vinculação de esquema mais recentes buscam eliminar elementos redundantes do esquema sem descartar informação relevante \cite{Yuan2025Knapsack,Wang2025LinkAlign} --- desafio análogo ao enfrentado pela etapa de compactação do esquema analítico do SIM (Seção 3.3).

\subsection{Recuperação Documental em Textos Governamentais}

A técnica RAG combina recuperação e geração textual \cite{Lewis2020RAG,Gao2024RAGSurvey,Asai2024SelfRAG,Sharma2025RAGSurvey}. Modelos densos como o DPR (\textit{Dense Passage Retrieval}) \cite{Karpukhin2020DPR} capturam similaridade semântica, mas exibem falhas sistemáticas em consultas contendo identificadores exatos, como números e anos de leis e decretos. Métodos híbridos que combinam representações densas e esparsas, a exemplo do BM25 \cite{Robertson2009BM25,Lin2021Hybrid,Lin2021Pyserini}, atenuam essa fragilidade. O SIM não propõe um novo algoritmo genérico de recuperação, mas avalia e adapta a busca híbrida combinando correspondência lexical exata para o cabeçalho do ato com similaridade vetorial para o teor normativo em Diários Oficiais municipais.

\subsection{GovTech e Posicionamento do SIM}

Iniciativas recentes de ciência de dados aplicada ao SICOM \cite{Lei2024SICOM} exploram auditoria e detecção de anomalias fiscais por especialistas. Sistemas como o GovChat abordam interfaces em linguagem natural para grafos de conhecimento governamentais \cite{KatsogiannisMeimaris2023GovChat}. O SIM diferencia-se dessas abordagens ao oferecer uma plataforma conversacional ponta-a-ponta voltada ao cidadão leigo, combinando dados estruturados contábeis e atos normativos não estruturados em uma arquitetura de orquestração leve, multi-instância e integrada a canais populares de mensageria.

\section{O SIM e a Instância Piloto de Caeté}

O SIM é projetado como uma plataforma modular e multi-instância para tecnologia cívica. A arquitetura foi concebida de forma genérica para operar sobre qualquer município brasileiro que disponibilize dados no padrão do SICOM ou esquemas contábeis correlatos. A cidade de Caeté (MG) foi selecionada como estudo de caso piloto por reunir acesso direto aos atores locais, base contábil completa de 2020 a 2025, acervo digitalizado de Diários Oficiais e lacunas documentadas de publicação de folha de pagamento e cadastro de servidores entre 2020 e 2022, que permitiram validar a resiliência do sistema em condições autênticas do setor público. As subseções a seguir descrevem a arquitetura operacional (Seção 3.1), as fontes de dados e o pipeline de ingestão (Seção 3.2) e as técnicas de engenharia de contexto empregadas (Seção 3.3).

\subsection{Arquitetura e Fluxo Operacional}

A arquitetura do SIM é organizada em três camadas funcionais, conforme ilustrado na Figura~\ref{fig:arquitetura}:

% Diagrama da arquitetura do SIM — incluído via \input
\begin{figure}[htbp]
    \centering
    \resizebox{0.92\linewidth}{!}{%
    \begin{tikzpicture}[
        font=\small,
        >={Stealth[length=2mm]},
        box/.style={draw, rounded corners=2pt, align=center, minimum width=2.5cm,
                    minimum height=1.0cm, fill=white, line width=0.5pt},
        store/.style={box, fill=black!7},
        flow/.style={->, thick, line width=0.6pt},
        biflow/.style={<->, thick, line width=0.6pt},
        elabel/.style={font=\footnotesize, fill=white, inner sep=1.5pt},
        zone/.style={rounded corners=5pt, draw=black!30, dash pattern=on 2pt off 2pt,
                     inner sep=0.3cm},
        zonelbl/.style={font=\scriptsize\bfseries, color=black!55, inner sep=2pt},
        node distance=1.05cm and 0.95cm
    ]
        % ===== Camada 1: canal conversacional (produção) =====
        \node[box, minimum width=2.6cm] (user) {Usuário\\\textit{WhatsApp}};
        \node[box, right=of user] (waha) {Gateway\\WAHA};
        \node[box, right=of waha]  (n8n)  {Automação\\n8n};
        \node[box, right=of n8n, minimum width=2.8cm] (orch)
              {Agente orquestrador\\[1pt]{\scriptsize\texttt{GPT-4o}}};

        \draw[flow] (user) -- (waha);
        \draw[flow] (waha) -- (n8n);
        \draw[flow] (n8n)  -- (orch);

        % retorno: resposta em linguagem natural (roteada acima da linha)
        \coordinate (top) at ($(orch.north)+(0,0.9)$);
        \draw[flow] (orch.north) -- (top) -| (user.north);
        \node[elabel] at ($(user.north)!0.5!(orch.north)+(0,0.9)$)
              {resposta em linguagem natural};

        % ===== Camada 2: recuperação e consulta =====
        \node[box, below=1.8cm of n8n,  minimum width=3.0cm] (rag)
              {RAG documental\\{\scriptsize Diário Oficial}};
        \node[box, below=1.8cm of orch, minimum width=3.0cm] (sql)
              {Text-to-SQL\\[1pt]{\scriptsize\texttt{GPT-4.1-mini}}};

        \coordinate (sp) at ($(orch.south)+(0,-0.75)$);
        \draw[flow] (orch.south) -- (sp);
        \draw[biflow] (sp) -| (rag.north);
        \draw[biflow] (sp) -- (sql.north);
        \node[elabel] at ([yshift=0.34cm]rag.north) {normativo};
        \node[elabel] at ([yshift=0.34cm]sql.north) {tabular / analítico};

        % ===== Camada 3: dados =====
        \node[store, below=0.9cm of rag] (qdrant)
              {Índice vetorial\\Qdrant};
        \node[store, below=0.9cm of sql] (duck)
              {API analítica\\DuckDB \textit{in-memory}};
        
        \node[store, minimum width=3.8cm, below=0.9cm of qdrant] (ingest)
              {Indexação Documental\\{\scriptsize Diário Oficial (PDFs / OCR)}};
        \node[store, minimum width=3.8cm, below=0.9cm of duck] (etl)
              {ETL SICOM e CGU\\{\scriptsize DDL e Tabelas Consolidadas}};

        \draw[biflow] (rag) -- (qdrant);
        \draw[biflow] (sql) -- (duck);
        \draw[flow] (ingest) -- (qdrant);
        \draw[flow] (etl) -- (duck);

        % ===== Zonas de fundo =====
        \begin{scope}[on background layer]
            \node[zone, fit=(user)(waha)(n8n)(orch)] (zA) {};
            \node[zone, fit=(rag)(sql)]               (zB) {};
            \node[zone, fit=(qdrant)(duck)(ingest)(etl)] (zC) {};
        \end{scope}
        \node[zonelbl, anchor=south west] at (zA.north west) {Canal conversacional (produção)};
        \node[zonelbl, anchor=south west] at (zB.north west) {Recuperação e consulta};
        \node[zonelbl, anchor=south west] at (zC.north west) {Camada de dados e ingestão};
    \end{tikzpicture}%
    }
    \caption{Arquitetura e fluxo operacional do SIM na instância piloto de Caeté.}
    \label{fig:arquitetura}
\end{figure}


\begin{enumerate}
    \item \textbf{Canal Conversacional (Produção):} O cidadão envia mensagens via WhatsApp. Um gateway de mensageria baseado em \texttt{WAHA} (\textit{WhatsApp HTTP API}) \cite{WAHA2024} repassa os eventos à plataforma de automação \texttt{n8n} \cite{n8n2024}, que aciona o agente orquestrador cognitivo implementado com \texttt{GPT-4o};
    \item \textbf{Recuperação e Consulta:} O orquestrador classifica a intenção do usuário em comunicação bidirecional com os subsistemas especializados. Perguntas normativas são roteadas ao subsistema RAG; consultas analíticas e quantitativas são encaminhadas ao subsistema Text-to-SQL baseado em \texttt{GPT-4.1-mini};
    \item \textbf{Camada de Dados e Ingestão:} O RAG opera sobre o índice vetorial \texttt{Qdrant} \cite{Qdrant2024}, alimentado pelo pipeline de indexação de Diários Oficiais. O subsistema Text-to-SQL executa consultas sobre o banco de dados analítico em memória \texttt{DuckDB} \cite{DuckDB2024}, abastecido pelo pipeline ETL.
\end{enumerate}

\subsection{Fontes de Dados e Pipeline ETL}

A instância Caeté cobre o período de 2020 a 2025, integrando: (i)~dados contábeis do SICOM referentes a receitas, despesas, empenhos, folha de pagamento, frota, licitações e contratos; (ii)~dados federais da CGU sobre transferências voluntárias e programas sociais, e (iii)~atos normativos do Diário Oficial.

Um pipeline ETL desenvolvido em Python consolida os arquivos tabulares, padroniza chaves de relacionamento, normaliza tipos de dados e gera o esquema analítico (\textit{Data Definition Language}, DDL). As lacunas de publicação de folha de pagamento e cadastro de servidores dos anos de 2020 a 2022 constituem limitações documentadas da base de dados municipal local.

Como exemplo de transformação, o pipeline realiza a higienização de chaves primárias ambíguas (remoção de pontuação e caracteres especiais em \texttt{cpf\_formatado}) e a conversão de datas numéricas no padrão \texttt{YYYYMMDD} para tipos relacionais \texttt{DATE} nativos do \texttt{DuckDB}, eliminando junções cartesianas indevidas em credores com registros heterogêneos. A Tabela~\ref{tab:registro_bruto} ilustra um registro bruto típico das tabelas SICOM, cuja nomenclatura verbosa e atributos de controle redundantes motivam a etapa de compactação do esquema analítico descrita na Seção 3.3.

\begin{table}[htbp]
\centering
\caption{Registro bruto da tabela \texttt{orgao.orgao} (SICOM, Caeté, 2023) antes da compactação do esquema}
\label{tab:registro_bruto}
\small
\begin{tabular}{ll}
\toprule
\textbf{Atributo (bruto)} & \textbf{Valor de exemplo} \\
\midrule
\texttt{seq\_orgao} & 304 \\
\texttt{num\_anoexercicio} & 2023 \\
\texttt{cod\_orgao} & 03 \\
\texttt{nom\_orgao} & SERVIÇO AUTÔNOMO DE ÁGUA E ESGOTO DE CAETÉ \\
\texttt{tipo\_orgao} & 03 - AUTARQUIA (EXCETO RPPS) \\
\texttt{cod\_municipio}* & 3110004 \\
\texttt{nom\_municipio}* & CAETÉ \\
\texttt{cod\_uf} & 31 \\
\texttt{sgl\_uf} & MG \\
\texttt{nom\_uf} & Minas Gerais \\
\texttt{dsc\_regiaoplanejamento} & CENTRAL \\
\bottomrule
\end{tabular}
\vspace{2pt}
\begin{minipage}{0.95\linewidth}
\footnotesize *Atributos removidos na compactação do esquema (Seção 3.3): a base é exclusiva do município de Caeté, portanto o código e o nome do município são redundantes em toda tabela que os contém --- assim como \texttt{cod\_uf}, \texttt{sgl\_uf} e \texttt{nom\_uf} nesta mesma tabela, que também não variam entre registros mas não são removidos pela regra de compactação por não fazerem parte da lista de atributos de controle identificada no pipeline ETL. A mesma poda remove ainda, em outras tabelas do esquema não ilustradas aqui, os atributos de controle de arquivo/lote \texttt{num\_versao\_arq} (tabelas de frota) e \texttt{num\_lote} (tabelas de licitação).
\end{minipage}
\end{table}


\subsection{Engenharia de Contexto}

Para garantir a geração precisa de SQL sem exceder o orçamento de tokens ou comprometer a latência síncrona do canal conversacional, foram avaliadas técnicas de engenharia de contexto:
\begin{itemize}
    \item \textbf{Compactação do Esquema Analítico:} Remoção de atributos redundantes de controle interno de arquivo/lote (por exemplo, \texttt{cod\_municipio}, \texttt{num\_versao\_arq}, \texttt{seq\_dat\_inclusao} e \texttt{num\_lote} --- metadados de extração sem valor analítico, já que cada instância opera sobre um único município) e condensação de metadados relacionais, preservando chaves primárias e tipos de dados essenciais para reduzir o consumo de tokens em aproximadamente 50\%;
    \item \textbf{Amostragem Inteligente (\textit{Smart Sampling}):} Inclusão de três a cinco exemplos reais de valores categóricos distintos por coluna relevante diretamente no cabeçalho do esquema, evitando alucinações de formato em cláusulas \texttt{WHERE};
    \item \textbf{Prompt A (Linha de Base Compacta):} Contém o esquema analítico compactado, os valores amostrados e um conjunto inicial de regras de domínio do padrão SICOM (conversão de datas \texttt{YYYYMMDD} com \texttt{CAST}, higiene de chaves e âncoras de \textit{join} setoriais), consumindo aproximadamente 16\,500 tokens;
    \item \textbf{Prompt B (Adaptação de Domínio):} Preserva a extensão do Prompt A ($\approx$16\,500 tokens) e estende suas regras de domínio com restrições determinísticas adicionais de precisão fina --- filtros de higiene condicionados ao nível de agregação da resposta, escopo de tabelas setoriais (como autarquias de saneamento SAAE e Restos a Pagar) e regras estritas de junção entre bases SICOM e CGU pelo identificador único \texttt{cpf\_formatado}. Esse refinamento incremental de regras de domínio é consistente com a tendência recorrente na literatura recente de Text-to-SQL, em que ganhos expressivos advêm de ajustes finos de adaptação de domínio e vinculação de esquema, e não apenas do aumento de escala do modelo subjacente \cite{Tian2025DomainAdapt,Hong2025Survey}.
\end{itemize}

\section{Metodologia de Validação Experimental}

Esta seção detalha o conjunto de referência (Seção 4.1), as métricas formais de avaliação (Seção 4.2) e a organização dos três lotes experimentais (Seção 4.3).

\subsection{Conjunto de Referência (\textit{Golden Dataset})}

O conjunto de referência Text-to-SQL reúne $N_{\text{SQL}} = 80$ casos executáveis estruturados a partir de 10 intenções analíticas centrais de relevância fiscal e administrativa (Tabela~\ref{tab:intencoes}), expandidas parametricamente com variações de anos, limites e filtros. A Tabela~\ref{tab:exemplos} apresenta exemplos representativos de consultas para cada nível de dificuldade e estrato.

\begin{table}[htbp]
\centering
\caption{Intenções analíticas centrais do conjunto Text-to-SQL ($N_{\text{SQL}} = 80$)}
\label{tab:intencoes}
\small
\begin{tabular}{clp{9cm}}
\toprule
\textbf{ID} & \textbf{Domínio Temático} & \textbf{Descrição da Consulta-Base} \\
\midrule
1 & Saúde Financeira & Cálculo do Índice de Saúde Financeira Municipal (ISF-M) \\
2 & Auditoria Social & Cruzamento de servidores públicos municipais com o Bolsa Família \\
3 & Compras Públicas & Cruzamento de fornecedores com maiores contratos SICOM e CGU \\
4 & Gestão de Frota & Agregação de gastos de manutenção e abastecimento por marca \\
5 & Educação & Cálculo do Índice de Eficiência em Educação (IEQ-C) \\
6 & Restos a Pagar & Consulta de restos a pagar cancelados e liquidados por exercício \\
7 & Licitações & Identificação das maiores licitações homologadas e vencedores \\
8 & Contratos & Identificação de contratos ativos com aditivos de acréscimo de valor \\
9 & Legislativo & Proporção de gastos da Câmara Municipal em relação à receita corrente \\
10 & Folha de Pagamento & Evolução temporal dos gastos consolidados com pessoal por secretaria \\
\bottomrule
\end{tabular}
\end{table}

\begin{table}[htbp]
\centering
\caption{Exemplos representativos de consultas do conjunto de referência}
\label{tab:exemplos}
\resizebox{\linewidth}{!}{%
\begin{tabular}{llp{7.5cm}l}
\toprule
\textbf{Subsistema} & \textbf{Estrato / Nível} & \textbf{Pergunta em Linguagem Natural} & \textbf{Alvo de Validação} \\
\midrule
Text-to-SQL & Fácil ($n=22$) & ``Qual foi a receita corrente total arrecadada em 2024?'' & Agregação simples em tabela única \\
Text-to-SQL & Médio ($n=44$) & ``Quais os 5 maiores credores da prefeitura em 2023?'' & \textit{Join} despesas $\times$ credores com \texttt{ORDER BY} \\
Text-to-SQL & Difícil ($n=14$) & ``Liste servidores ativos que receberam benefício social em 2023.'' & \textit{Join} complexo SICOM $\times$ CGU por CPF \\
\midrule
RAG & \texttt{com\_numero} ($n=32$) & ``Qual o teor do Decreto Municipal nº 142 de 2023?'' & Recuperação exata por número e ano \\
RAG & \texttt{tematico} ($n=8$) & ``Quais as normas publicadas sobre concurso público em 2024?'' & Similaridade semântica sobre normas \\
\bottomrule
\end{tabular}%
}
\end{table}

Para o subsistema RAG, construiu-se conjunto ampliado com $N_{\text{RAG}} = 40$ perguntas sobre os $9\,346$ trechos (\textit{chunks}) do Diário Oficial indexados no \texttt{Qdrant}: (i)~\texttt{com\_numero} ($n = 32$), composto por decretos sorteados uniformemente do índice vetorial, e (ii)~\texttt{tematico} ($n = 8$), contendo perguntas curadas sobre tópicos normativos sem menção a número de ato.

\subsection{Métricas de Avaliação}

Seja $\text{SQL}_i^{(g)}$ a consulta gerada pelo modelo e $\text{SQL}_i^{(r)}$ a consulta de referência para a pergunta $i \in \{1,\dots,N_{\text{SQL}}\}$. Definem-se:

\textbf{Taxa de SQL Válida (VSR):} Proporção de consultas geradas que executam sem erro sintático no motor analítico:
\begin{equation}
\text{VSR} = \frac{1}{N_{\text{SQL}}} \sum_{i=1}^{N_{\text{SQL}}} \mathbb{I}\!\left[\text{SQL}_i^{(g)} \text{ executa sem erro}\right]
\end{equation}
onde $\mathbb{I}[\cdot]$ denota a função indicadora, assumindo valor 1 se a condição for verdadeira e 0 caso contrário.

\textbf{Equivalência Funcional (EF):} Sejam $R_i$ o resultado tabular retornado pela consulta de referência (gabarito) para a pergunta $i$, e $G_i$ o resultado tabular retornado pela consulta gerada pelo modelo. Seja ainda $C_i$ o conjunto de colunas de resposta declaradas como relevantes no gabarito da pergunta $i$ --- ou seja, o subconjunto de colunas de $R_i$ que efetivamente respondem à pergunta formulada, excluindo colunas auxiliares de ordenação ou depuração eventualmente presentes na consulta de referência. A EF verifica se $G_i$ entrega, para cada coluna em $C_i$, o mesmo conteúdo informacional presente em $R_i$. Diferentemente do EX estrito utilizado em Spider e BIRD --- que exige nomes de colunas idênticos ao gabarito ---, a EF relaxa essa exigência em três pontos: (i) \textbf{correspondência posicional}, isto é, cada coluna declarada em $C_i$ é localizada em $G_i$ pela sua posição ordinal (a $k$-ésima coluna de $C_i$ é comparada à $k$-ésima coluna de $G_i$), e não pelo nome ou \textit{alias} atribuído pelo modelo; (ii) \textbf{cardinalidade de linhas compatível} ($|R_i| = |G_i|$), sem exigir que as linhas estejam na mesma ordem de apresentação, e (iii) \textbf{tolerância numérica relativa} em valores monetários, definida como $\tau = 0{,}5\%$, para absorver arredondamentos de ponto flutuante inerentes a operações contábeis. Após normalização celular $\nu(\cdot)$ (conversão de tipos e remoção de formatação de string), a igualdade aproximada entre um valor gerado $a$ e o valor de referência correspondente $b$ é definida por:
\begin{equation}
\frac{|\nu(a) - \nu(b)|}{\max(|\nu(b)|, \varepsilon)} \leq \tau, \quad \text{EF} = \frac{1}{N_{\text{SQL}}} \sum_{i=1}^{N_{\text{SQL}}} \mathbb{I}\!\left[\text{EF}_i = 1\right]
\end{equation}
em que $\varepsilon = 10^{-9}$ é uma constante de estabilização numérica, sem significado fiscal, que evita divisão por zero quando o valor de referência $\nu(b)$ é nulo, e $\tau = 0{,}005$ é a tolerância relativa de 0,5\% já definida --- os dois símbolos desempenham papéis distintos e não devem ser confundidos: $\varepsilon$ protege o cálculo, $\tau$ define o critério de aceitação. Define-se $\text{EF}_i = 1$ quando, para toda coluna em $C_i$, o valor correspondente em $G_i$ satisfaz essa desigualdade. A verificação da EF é integralmente automática, por comparação programática entre os vetores de coluna do resultado gerado e do gabarito (implementada como comparação de vetores ordenados por coluna, não linha a linha), sem julgamento humano na avaliação individual dos casos.

\textbf{Métricas RAG:} Para o subsistema de recuperação documental, avaliam-se o \textbf{Recall@}$k$ (presença do trecho esperado entre os $k$ primeiros resultados retornados, com $k=5$) e o \textbf{MRR} (\textit{Mean Reciprocal Rank}, inverso da posição do primeiro acerto no ranking) \cite{BaezaYates2011}.

\textbf{Avaliação por Juiz Automático (JAR):} A taxa de concordância do juiz (\textit{Judge Agreement Rate}, JAR) afere a fidelidade da resposta final em linguagem natural produzida pelo orquestrador, avaliada pelo \texttt{GPT-4o} atuando como juiz automático sobre a amostra estratificada. Para cada caso, o juiz produz duas saídas independentes: (i) notas em escala Likert de 1 a 5 em quatro dimensões --- adequação semântica, fidelidade aos dados tabulares e normativos, completude da informação e clareza da linguagem para o cidadão leigo; e (ii) um veredito categórico (\textit{correto}, \textit{parcialmente correto} ou \textit{incorreto}). O JAR é definido como a proporção de casos cujo veredito é \textit{correto} ou \textit{parcialmente correto} --- ou seja, a taxa de respostas que o juiz considera aceitáveis frente ao padrão esperado de resposta correta, e não uma medida de concordância entre dois avaliadores distintos. As notas Likert são reportadas separadamente, como indicador complementar de qualidade, e não participam do cálculo do JAR \cite{Zheng2023LLMJudge}.

Todas as proporções são reportadas com intervalos de confiança (IC) de Wilson a 95\% \cite{Wilson1927,Agresti1998}. Comparações pareadas entre variantes utilizam o teste de McNemar \cite{Agresti1998}.

\subsection{Organização dos Lotes Experimentais}

A avaliação estruturou-se em três lotes (Tabela~\ref{tab:lotes}), isolando variáveis experimentais controladas.

\begin{table}[htbp]
\centering
\caption{Organização dos lotes experimentais}
\label{tab:lotes}
\small
\begin{tabular}{cllp{7.2cm}}
\toprule
\textbf{Lote} & \textbf{Variável Isolada} & \textbf{Domínio} & \textbf{Configurações Avaliadas} \\
\midrule
I & Modelo e contexto & Text-to-SQL & E1 (GPT-4.1-mini, compacto); E2 (estendido); E3 (4o-mini); E4 (4o) \\
II & Engenharia de \textit{prompt} & Text-to-SQL & Prompt A (Baseline E1) \textit{versus} Prompt B (regras de domínio) \\
III & Estratégia de recuperação & RAG & Densa; deduplicação; filtro por tipo; busca híbrida \\
\bottomrule
\end{tabular}
\end{table}

\section{Resultados Experimentais e Análise}

Esta seção apresenta os resultados quantitativos obtidos nos lotes experimentais de Text-to-SQL (Seções 5.1 e 5.2), recuperação documental (Seção 5.3) e validação da orquestração conversacional (Seção 5.4).

\subsection{Lote I --- Comparação de Modelo e Contexto (E1--E4)}

A Tabela~\ref{tab:lote1} apresenta os resultados do Lote I sobre os 80 casos do conjunto Text-to-SQL.

\begin{table}[htbp]
\centering
\caption{Lote I --- Comparativo de modelos e contextos ($N_{\text{SQL}} = 80$)}
\label{tab:lote1}
\resizebox{\linewidth}{!}{%
\begin{tabular}{lrrrrr}
\toprule
\textbf{Experimento / Configuração} & \textbf{VSR (\%)} & \textbf{IC 95\% VSR} & \textbf{EF (\%)} & \textbf{IC 95\% EF} & \textbf{Latência (s)} \\
\midrule
E1 --- Baseline (GPT-4.1-mini, Prompt A) & 97,50 & 91,3--99,3 & 62,50 & 51,5--72,3 & 6,2 \\
E2 --- Contexto estendido (GPT-4.1-mini) & 96,25 & 89,5--98,7 & 58,75 & 47,8--68,9 & 11,4 \\
E3 --- Modelo leve (GPT-4o-mini) & 90,00 & 81,5--94,8 & 46,25 & 35,7--57,1 & 3,8 \\
E4 --- Modelo avançado (GPT-4o) & 97,50 & 91,3--99,3 & 57,50 & 46,6--67,7 & 29,5 \\
\bottomrule
\end{tabular}%
}
\end{table}

As quatro configurações do Lote I isolam a variável de modelo e contexto mantendo o restante do protocolo fixo. E1, E3 e E4 utilizam exatamente o mesmo Prompt A descrito na Seção 3.3 ($\approx$16\,500 tokens), variando apenas o modelo subjacente (\texttt{GPT-4.1-mini}, \texttt{GPT-4o-mini} e \texttt{GPT-4o}, respectivamente); E2 substitui o Prompt A por uma variante estendida própria, com $\approx$30\,000 tokens, que preserva o mesmo esquema e as mesmas regras de domínio, mas inclui documentação adicional de colunas e exemplos amostrados em maior quantidade por tabela.

O baseline E1 estabeleceu desempenho sólido com VSR de 97,5\% e EF de 62,5\%. O contexto estendido (E2, $\approx$30k tokens, prompt próprio conforme descrito acima) apresentou EF pontual de 58,75\%; embora a sobreposição dos intervalos de confiança (51,5--72,3\% em E1 vs. 47,8--68,9\% em E2) não permita afirmar superioridade estatística estrita, a tendência observada é coerente com o fenômeno de dispersão de atenção em contextos excessivamente longos (\textit{lost in the middle}) \cite{Liu2023LostMiddle}. O \texttt{GPT-4o-mini} (E3, mesmo Prompt A de E1) mostrou-se insuficiente para a complexidade do esquema SICOM (EF de 46,25\%). O \texttt{GPT-4o} (E4, mesmo Prompt A de E1) não superou o modelo compacto em EF (57,50\%) e apresentou latência média de geração de 29,5 segundos --- inviável para operação síncrona no WhatsApp. Esses resultados fixaram o \texttt{GPT-4.1-mini} como modelo de referência para os experimentos subsequentes.

A Tabela~\ref{tab:diff} detalha o desempenho do baseline E1 por nível de dificuldade. Observa-se que a degradação de desempenho é gradual, concentrando-se em consultas de cruzamento entre tabelas e agregações temporais complexas.

\begin{table}[htbp]
\centering
\caption{E1 --- Desempenho por nível de dificuldade ($N_{\text{SQL}} = 80$)}
\label{tab:diff}
\small
\begin{tabular}{lrrrr}
\toprule
\textbf{Dificuldade} & \textbf{Casos ($n$)} & \textbf{VSR (\%)} & \textbf{EF (\%)} & \textbf{IC 95\% EF} \\
\midrule
Fácil & 22 & 100,00 & 68,18 & 47,3--83,6 \\
Médio & 44 & 97,73 & 61,36 & 46,6--74,3 \\
Difícil & 14 & 92,86 & 57,14 & 32,6--78,6 \\
\bottomrule
\end{tabular}
\end{table}

\subsection{Lote II --- Engenharia de Contexto e Adaptação de Domínio}

Mantendo fixos o modelo \texttt{GPT-4.1-mini} e o conjunto de teste ($N_{\text{SQL}} = 80$), o Lote II avaliou o impacto das regras de domínio do Prompt B frente ao Prompt A, conforme apresentado na Tabela~\ref{tab:lote2}.

\begin{table}[htbp]
\centering
\caption{Lote II --- Comparativo de engenharia de \textit{prompt} ($N_{\text{SQL}} = 80$)}
\label{tab:lote2}
\small
\begin{tabular}{lrrrr}
\toprule
\textbf{Configuração de Prompt} & \textbf{VSR (\%)} & \textbf{IC 95\% VSR} & \textbf{EF (\%)} & \textbf{IC 95\% EF} \\
\midrule
Prompt A --- Baseline compacto (E1) & 97,50 & 91,3--99,3 & 62,50 & 51,5--72,3 \\
Prompt B --- Regras de domínio (Final) & 100,00 & 95,4--100,0 & 72,50 & 61,9--81,1 \\
\bottomrule
\end{tabular}
\end{table}

O Prompt B alcançou 100\% de SQLs válidas e elevou a equivalência funcional em 10 pontos percentuais (72,50\%), sem acréscimo de latência. Na análise pareada de pares discordantes, registraram-se 12 consultas corrigidas pelo Prompt B contra 5 regressões (teste de McNemar: $p = 0{,}146$).

Ressalta-se que o desenvolvimento do Prompt B não representou um ajuste empírico restrito (\textit{overfitting}) às perguntas da amostra. Pelo contrário, atuou como uma \textbf{adaptação de domínio contábil}, incorporando restrições ontológicas e convenções normativas do padrão SICOM/TCE-MG compartilhadas por todos os municípios mineiros (Tabela~\ref{tab:erros}).

\begin{table}[htbp]
\centering
\caption{Taxonomia de falhas do baseline E1 sanadas pelo Prompt B}
\label{tab:erros}
\resizebox{\linewidth}{!}{%
\begin{tabular}{p{1.6cm}p{4.0cm}p{4.2cm}p{4.2cm}}
\toprule
\textbf{Caso / ID} & \textbf{Consulta (Resumo)} & \textbf{Falha no Baseline E1} & \textbf{Regra de Domínio (Prompt B)} \\
\midrule
ID 8 & Contratos com aditivo de acréscimo & Comparação de \texttt{BIGINT} com \texttt{DATE} em datas YYYYMMDD & Conversão explícita com \texttt{CAST} de datas SICOM \\
ID 2 & Servidores ativos $\times$ Bolsa Família & Cardinalidade divergente por homônimos em \textit{join} por nome & Junção por chave única documentada (\texttt{cpf\_formatado}) \\
ID 7 & Maiores licitações homologadas & Agregação incorreta sem \textit{join} com itens de licitação & Âncoras de \textit{join} em tabelas setoriais \\
ID 10 & Evolução mensal da folha & Filtro de higiene indevido restringindo universo & Escopo condicional de filtros ao formato de agregação \\
\bottomrule
\end{tabular}%
}
\end{table}

\subsection{Lote III --- Recuperação Documental sobre o Diário Oficial}

O Lote III comparou quatro estratégias de recuperação sobre o conjunto de referência RAG ampliado ($N_{\text{RAG}} = 40$), mantendo o índice vetorial e o modelo de \textit{embedding} constantes (Tabela~\ref{tab:lote3}).

\begin{table}[htbp]
\centering
\caption{Lote III --- Comparativo de estratégias de recuperação RAG ($N_{\text{RAG}} = 40$, $k = 5$)}
\label{tab:lote3}
\small
\begin{tabular}{lrrr}
\toprule
\textbf{Estratégia de Recuperação} & \textbf{Recall@5 (\%)} & \textbf{IC 95\% Recall} & \textbf{MRR (\%)} \\
\midrule
Busca densa isolada & 22,50 & 12,3--37,5 & 13,17 \\
Deduplicação de \textit{chunks} idênticos & 22,50 & 12,3--37,5 & 14,79 \\
Filtro determinístico por tipo de ato + dedup & 22,50 & 12,3--37,5 & 15,42 \\
\textbf{Busca híbrida por nº de ato + dedup + filtro} & \textbf{100,00} & \textbf{91,2--100,0} & \textbf{97,50} \\
\bottomrule
\end{tabular}
\end{table}

A desagregação por estrato explicita o fenômeno: no estrato \texttt{com\_numero} ($n = 32$), a busca densa pura obteve apenas 3,1\% de Recall@5 (1/32). Modelos de \textit{embedding} falham em preservar a correspondência exata de identificadores numéricos, retornando decretos com teor similar, porém numeração incorreta. A busca híbrida (extração do identificador e priorização do cabeçalho do ato) elevou o Recall@5 a 100\% e o MRR a 97,5\%. No estrato \texttt{tematico} ($n = 8$), a busca densa já alcançava 100\% de Recall@5, sendo refinada pela deduplicação.

\subsection{Validação Qualitativa da Orquestração e Resposta Natural}

Esta seção avalia uma etapa do pipeline distinta e posterior àquela mensurada pela EF: enquanto a EF (Seção 4.2) verifica exclusivamente a corretude do resultado tabular bruto retornado pelo subsistema Text-to-SQL, de forma automática e determinística, a avaliação aqui descrita mede a fidelidade da \emph{resposta final em linguagem natural} produzida pelo orquestrador cognitivo a partir desse resultado tabular. As duas avaliações são necessariamente distintas e complementares: é possível gerar uma consulta SQL funcionalmente correta ($\text{EF}_i = 1$) e, ainda assim, produzir uma resposta textual inadequada ao interpretar, resumir ou formatar esse resultado para o cidadão --- falha que a EF, por medir apenas o resultado tabular, não é capaz de capturar. A taxa de concordância do juiz (JAR), definida na Seção 4.2, existe justamente para cobrir essa lacuna.

A avaliação por juiz automático (\texttt{GPT-4o}) sobre a amostra estratificada do E1 ($n = 30$; dez casos sorteados dentre os já concluídos em cada um dos três níveis de dificuldade --- fácil, médio e difícil --- que estruturam a Tabela~\ref{tab:diff}) registrou a tripla: pergunta $\rightarrow$ resultado \texttt{DuckDB} $\rightarrow$ resposta textual gerada. Obteve-se taxa de concordância (JAR) de 100\% (30/30; IC 95\%: 88,6--100\%), com notas médias de 4,87/5 em adequação semântica, 4,87/5 em fidelidade aos dados tabulares, 4,73/5 em completude da informação e 4,77/5 em clareza da linguagem para o cidadão \cite{Zheng2023LLMJudge}. Para o subsistema RAG, a avaliação qualitativa sobre o conjunto de referência completo ($N_{\text{RAG}} = 40$, derivado de \texttt{golden\_rag\_v1.3.csv}) registrou 100\% de Recall@5 e MRR de 97,5\% na busca híbrida, comprovando a ausência de alucinações e a fidelidade textual aos trechos recuperados do Diário Oficial.

\section{Discussão}

\subsection{Contextualização em Relação a Benchmarks Acadêmicos}

A equivalência funcional de 72,5\% alcançada pelo SIM situa-se na mesma ordem de grandeza dos resultados reportados em \textit{benchmarks} internacionais complexos como o BIRD, onde modelos comerciais atingem EX entre 50\% e 55\% \cite{Li2024BIRD}, situando-se na mesma ordem de grandeza reportada em trabalhos recentes de 2026 sobre interfaces no setor público \cite{Zheng2026GovLLM}. Ressalta-se que a EF não é diretamente comparável ao EX estrito: no baseline E1 (Tabela~\ref{tab:lote1}), a EX estrita foi de apenas 11,25\% (9 de 80 casos), pois o modelo frequentemente nomeia colunas com aliases descritivos diferentes do gabarito, embora os valores retornados fossem idênticos aos demandados pelo cidadão. A EF mostrou-se a métrica apropriada para mensurar a utilidade prática em tecnologia cívica.

\subsection{Achados Empíricos Centrais}

Os experimentos consolidam dois achados principais:
\begin{enumerate}
    \item \textbf{Engenharia de contexto supera troca de modelo:} Inserir convenções de negócio e regras do SICOM no \textit{prompt} recuperou +10\,pp de EF sobre o \texttt{GPT-4.1-mini} (Prompt A: 62,50\% $\rightarrow$ Prompt B: 72,50\%, Tabela~\ref{tab:lote2}), superando em EF o \texttt{GPT-4o} sem adaptação de domínio (E4: 57,50\%, Tabela~\ref{tab:lote1}) e reduzindo o tempo de resposta em aproximadamente 79\% frente a esse mesmo modelo (de 29,5\,s em E4 para 6,2\,s, latência de referência do Prompt A/E1 preservada pelo Prompt B);
    \item \textbf{Necessidade estrita de busca híbrida em Diários Oficiais:} O colapso da busca densa pura em atos numerados (22,5\% Recall@5 agregado e 3,1\% no estrato com número) evidencia que sistemas GovTech para documentos oficiais exigem indexação híbrida estruturada.
\end{enumerate}

\subsection{Confiabilidade, Riscos e Segurança}

Em canais públicos como o WhatsApp, a taxa de 27,5\% de falha funcional residual (em consultas difíceis e cruzamentos complexos) impõe riscos de desinformação fiscal. Como mitigação de engenharia, recomenda-se: (i)~exibição opcional da consulta SQL gerada para auditoria cidadã; (ii)~declaração de incerteza em cruzamentos com cardinalidade ambígua, e (iii)~mecanismos de proteção contra injeção de instruções (\textit{prompt injection}) na camada do orquestrador.

\subsection{Ameaças à Validade e Limitações}

Reconhecem-se como limitações: (i)~conjunto de referência derivado de 10 intenções analíticas em uma única instância municipal (Caeté); (ii)~o ganho do Prompt B foi aferido no mesmo conjunto de referência usado na taxonomia de falhas, e o tamanho amostral ($N_{\text{SQL}} = 80$) resulta em poder estatístico limitado para a comparação pareada, que não atingiu significância estatística ($p = 0{,}146$), e (iii)~a validação em produção no WhatsApp possui caráter qualitativo de engenharia, não tendo sido coletadas métricas de adoção populacional em larga escala.

\section{Conclusão}

Este artigo apresentou o SIM (Sistema de Informações Municipais), uma plataforma conversacional multi-instância que integra recuperação documental sobre Diários Oficiais e consultas analíticas Text-to-SQL sobre dados contábeis abertos (SICOM e CGU). A solução visa reduzir a barreira técnica de acesso à informação pública para o cidadão comum, operando via WhatsApp com orquestração cognitiva por modelos de linguagem.

A validação experimental em três lotes comprovou a viabilidade funcional da arquitetura. No subsistema Text-to-SQL, demonstrou-se que a engenharia de contexto especializada com adaptação de domínio eleva a equivalência funcional para 72,5\% com 100\% de SQLs válidas, superando modelos proprietários maiores com substancial economia de latência e custo. No subsistema RAG, evidenciou-se a falha crítica da busca densa isolada para atos governamentais numerados (Recall@5 de 3,1\%), solucionada pela adoção de recuperação híbrida (Recall@5 de 100\% e MRR de 97,5\%). O sistema encontra-se implantado em operação no canal WhatsApp da instância piloto de Caeté, confirmando a estabilidade da orquestração cognitiva em ambiente real.

Como trabalhos futuros, destacam-se: (i)~a realização de estudos de experiência do usuário e adoção populacional com cidadãos de Caeté; (ii)~a extensão do SIM para outras prefeituras mineiras sob o mesmo padrão SICOM, e (iii)~a incorporação de mecanismos automáticos de verificação formal de integridade fiscal antes do envio da resposta textual ao cidadão.

\section*{Disponibilidade de Dados e Código}

O código-fonte completo da plataforma SIM, as rotinas de ingestão ETL, os conjuntos de referência (\textit{golden datasets}) Text-to-SQL e RAG, e os scripts para replicação offline dos experimentos estão disponíveis publicamente sob licença aberta em: \url{https://github.com/rafa-rez/tcc-sistema-de-informa-es-municipais}.

\begingroup
\footnotesize
\bibliography{../referencias/referencias_sim}
\endgroup

\enddocument
